\section{Ethics and Privacy}\label{sec:ethics}

\paragraph{Bystanders.} A head-worn camera sees people who did not consent to being recorded, and bystanders' concerns about camera glasses are well documented~\cite{denning2014situ,hoyle2014privacy}. SpatialNet runs on the glasses, so localisation needs no frame to leave the device, and a later step can be limited to the target crop and to moments when the wearer points. SpatialNet is trained to ignore other people's hands; it does not detect, identify or track people. On-device processing does not remove the need for visible recording indicators; the test unit has user- and bystander-facing LEDs.

\paragraph{Data.} EgoPoint is used under the access terms of its authors' release, for research only. We do not redistribute its images or our converted labels; our conversion script regenerates the labels from the original release. The images shown in this paper come from its test and training splits. \tbd{Real-frame check (Section~\ref{sec:realframes}): state who recorded the frames and how they consented, the ethics approval or exemption, how long the frames are kept, and how bystanders were kept out of view or blurred.}

\paragraph{Availability.} We will release the code, the dataset conversion script, the trained weights and the per-image evaluation outputs, from which every table can be recomputed, upon publication.

\section{Conclusion}\label{sec:conclusion}
SpatialNet localises the object a smart-glasses wearer points at from a single RGB frame, with one forward pass of a 0.72M-parameter CNN and no prompt, language model or external hand model, and it runs on a pre-production pair of glasses in 60--80\,ms per frame. It is far less accurate than large GPU models, and compact detectors of the same compute draw tighter boxes. But it picks the labelled object more often than those detectors, a recogniser names the object from its crops as often as from theirs, and it returns the hand, the fingertip and the direction from the same pass. At this scale, input resolution is the largest lever, instance anchoring sharpens the hand box and the direction, training with distractor hands carries over to other models, and an explicit ray prior is not needed. At least 41\% of the EgoPoint test images are COCO photographs, which benchmarks on this dataset should take into account. The main open problems are pointing scenes that EgoPoint does not contain, real bystander hands, and evidence that on-device pointing helps people in daily use.
